% vim: set ft=tex tabstop=4 shiftwidth=4 noexpandtab:

\input{../../include/latex/tikz/header.tex}

% opening %{{{1

\begin{document}
\begin{tikzpicture}[scale=1.0]

% parameters %{{{1

	\def\ang{200}
	\def\ngl{-30}
	\def\gle{120}

	\def\radius{3.5}

	\def\lenBisector{6.0cm}

	%\tikzmath{
		%\var = 2.1;
	%}

% coordinates %{{{1

	\coordinate (A) at (\ang:\radius);
	\coordinate (B) at (\ngl:\radius);
	\coordinate (C) at (\gle:\radius);

	\coordinate (lenAB) at ($(A)!\lenBisector!(B)$);
	\coordinate (lenAC) at ($(A)!\lenBisector!(C)$);
	\coordinate (XbisectorBAC) at ($(lenAB)!0.5!(lenAC)$);

	\coordinate (lenBA) at ($(B)!\lenBisector!(A)$);
	\coordinate (lenBC) at ($(B)!\lenBisector!(C)$);
	\coordinate (XbisectorCBA) at ($(lenBA)!0.5!(lenBC)$);

	\coordinate (lenCA) at ($(C)!\lenBisector!(A)$);
	\coordinate (lenCB) at ($(C)!\lenBisector!(B)$);
	\coordinate (XbisectorACB) at ($(lenCA)!0.5!(lenCB)$);

% intersections %{{{1

	\begin{pgfinterruptboundingbox}
		\path[name path=line1] (A) -- (XbisectorBAC);
		\path[name path=line2] (B) -- (XbisectorCBA);
		\path[name intersections={of=line1 and line2, by=O}];
	\end{pgfinterruptboundingbox}

	% ---- projection of O onto triangle sides are
	% ---- intersections of inscribed circle and sides
	% ---- also to compute the inscribed circle radius

	\coordinate (P) at ($ (A)!(O)!(B) $);
	\coordinate (S) at ($ (B)!(O)!(C) $);
	\coordinate (T) at ($ (C)!(O)!(A) $);

% distance %{{{1

	\path let \p1 = ($ (P) - (O) $)
	in \pgfextra{
		\pgfmathsetmacro{\radiusInscribed}{veclen(\x1,\y1)/1cm}
		\global\let\radiusInscribed\radiusInscribed
	};

% points, dots, vertices %{{{1

	\fill (A) circle (0.4mm);
	\fill (B) circle (0.4mm);
	\fill (C) circle (0.4mm);

	%\fill (XbisectorBAC) circle (0.4mm);
	%\fill (XbisectorCBA) circle (0.4mm);
	%\fill (XbisectorACB) circle (0.4mm);

	\fill (O) circle (0.4mm);

	\fill (P) circle (0.4mm);
	\fill (S) circle (0.4mm);
	\fill (T) circle (0.4mm);

% segments, sides, lines %{{{1

	\draw (A) -- (B) -- (C) -- cycle;
	\draw[dashed] (O) -- (P);
	\draw[dashed] (O) -- (S);
	\draw[dashed] (O) -- (T);

% circles %{{{1

	\path[draw] let \p1 = ($ (O) - (T) $)
	in (O) circle ({veclen(\x1,\y1)});

% circular arcs %{{{2

	%\draw (A) ++(\ang:\radius) arc (\ang:\gle:\radius);

% points, dots, vertices labels %{{{1

	\node[below] at (A) {$A$};
	\node[below] at (B) {$B$};
	\node[above] at (C) {$C$};

	\node[below left] at (O) {$O$};

	\node[below] at (P) {$P$};
	\node[above right] at (S) {$S$};
	\node[above left] at (T) {$T$};

% segments, sides, lines labels %{{{1

	%\node[above] at ($ (A)!0.5!(B) $) {$d$};

% segments, sides, lines marks %{{{1

	%\tkzMarkSegments[mark=||, size=2pt](A,B)
	%\tkzMarkSegments[mark=||, size=2pt](C,D)

% circles labels %{{{1

	\node at ($ (O) + (-30:\radiusInscribed + 0.3) $) {$\mathscr{C}$};

% circular arcs labels %{{{2

	%\node at ($ (O) + (45:\radius + 0.3) $) {$\mathscr{C}$};

% angles labels %{{{1

	%\pic[draw, ->, "$\alpha$", angle radius=0.8cm, angle eccentricity=1.5]
	%{angle = B--A--C};

% right angles markers %{{{2

	%\pic [draw, angle radius=7pt, angle eccentricity=1]
	%{right angle = B--P--O};

	%\pic [draw, angle radius=7pt, angle eccentricity=1]
	%{right angle = C--S--O};

	%\pic [draw, angle radius=7pt, angle eccentricity=1]
	%{right angle = A--T--O};

% closing %{{{1

\end{tikzpicture}
\end{document}
